% =====================================================================
%  CHƯƠNG 4 — PHÁT HIỆN XE VÀ BIỂN SỐ          | Phụ trách: Nhật
%  Ngân sách: mục tiêu 5 trang.
%  Nguồn số nội bộ (không in ra báo cáo):
%    experiments.csv (test split, 2026-08-13); results.csv 6 run 50 epoch;
%    a1_dedup_eval.csv (2026-09-12); phash_report.txt (230 cặp Hamming<=5).
% =====================================================================
\chapter{Phát Hiện Xe và Biển Số}
\label{ch:phathien}

Chương này so sánh YOLOv8n và YOLO26n cho phát hiện biển số trên bộ A1, đánh giá trên cả
tập kiểm tra gốc lẫn tập đã khử trùng lặp gần.

\section{Kiến Trúc Mô Hình}
\label{sec:phathien-kientruc}

Kiến trúc hai mô hình đã trình bày ở Mục~\ref{sec:yolov8} và Mục~\ref{sec:yolo26}. Điểm
liên quan trực tiếp đến bài toán là YOLOv8n gán nhãn dương bằng Task-Aligned
Assigner~\autocite{jocher2023yolov8,feng2021tood}, còn YOLO26n bỏ NMS, bỏ DFL và dùng STAL ưu
tiên vật thể nhỏ~\autocite{jocher2026yolo26,chakrabarty2026yolo26nmsfree}, phù hợp với biển số
vốn nhỏ so với khung ảnh~\autocite{zhu2025licenseplate}. Cả hai dùng biến thể nano để vừa
ngân sách độ trễ trên Raspberry Pi 5 (Mục~\ref{sec:pi5bench}).

\section{Thiết Lập Huấn Luyện}
\label{sec:phathien-huanluyen}

Đề tài huấn luyện trên A1 (4.578 ảnh, 5.200 đối tượng biển số) với hai lớp là biển một hàng
(đặc trưng ô tô) và biển hai hàng (đặc trưng xe máy), chia 3.433 ảnh huấn luyện, 573 kiểm
định, 572 kiểm tra. Bảng~\ref{tab:det-hyperparam} liệt kê các siêu tham số chính; mô hình
khởi tạo từ trọng số pretrained COCO, mỗi kiến trúc huấn luyện với ba seed (0, 1, 2) để đánh
giá tính ổn định. Chi tiết augmentation ở Phụ lục~\ref{app:huanluyen}.

\begin{table}[htbp]
\centering
\caption{Siêu tham số huấn luyện detector biển số.}

\label{tab:det-hyperparam}
\begin{tabular}{ll}
\toprule
\textbf{Tham số} & \textbf{Giá trị} \\
\midrule
Kích thước ảnh đầu vào & 640 \\
Số epoch & 50 \\
Batch size & 16 \\
Patience (early stopping) & 10 \\
Optimizer & SGD (mặc định Ultralytics) \\
Learning rate khởi đầu & 0,01 \\
Learning rate cuối & 0,0001 \\
Momentum & 0,937 \\
Weight decay & 0,0005 \\
Warmup epochs & 3 \\
\bottomrule
\end{tabular}
\end{table}

\section{Độ Đo Đánh Giá}
\label{sec:phathien-dodo}

Precision $P = TP/(TP+FP)$ đo tỉ lệ dự đoán đúng trong các dự đoán dương; Recall
$R = TP/(TP+FN)$ đo tỉ lệ biển thật được phát hiện. \textbf{mAP@0,5} là diện tích dưới đường cong Precision-Recall tại ngưỡng IoU = 0,5,
trung bình trên các lớp.
\textbf{mAP@0,5:0,95} lấy trung bình mAP tại các ngưỡng IoU từ 0,5 đến 0,95 bước 0,05, phạt
nặng hơn các bounding box lỏng lẻo nên phản ánh chất lượng định vị chặt hơn. Giá trị mỗi
epoch tính trên tập kiểm định; kết quả cuối cùng báo cáo trên tập kiểm tra.

\section{Kết Quả Phát Hiện}
\label{sec:phathien-ketqua}

\subsection{So Sánh Hai Mô Hình trên Ba Seed}

\begin{table}[htbp]
\centering
\caption{Kết quả phát hiện biển số trên tập kiểm định (epoch 50) theo từng seed.}

\label{tab:det-results}
\begin{tabular}{llcccc}
\toprule
\textbf{Mô hình} & \textbf{Seed} & \textbf{P} & \textbf{R} & \textbf{mAP@0,5} & \textbf{mAP@0,5:0,95} \\
\midrule
\multirow{4}{*}{YOLOv8n} & 0   & 0,988 & 0,987 & 0,994 & 0,867 \\
                          & 1   & 0,990 & 0,992 & 0,994 & 0,873 \\
                          & 2   & 0,986 & 0,991 & 0,994 & 0,866 \\
                          & \textit{Trung bình} & \textit{0,988} & \textit{0,990} & \textit{0,994} & \textit{0,869} \\
\midrule
\multirow{4}{*}{YOLO26n}  & 0   & 0,982 & 0,969 & 0,993 & 0,879 \\
                           & 1   & 0,974 & 0,977 & 0,992 & 0,880 \\
                           & 2   & 0,969 & 0,988 & 0,993 & 0,884 \\
                           & \textit{Trung bình} & \textit{0,975} & \textit{0,978} & \textit{0,993} & \textit{0,881} \\
\bottomrule
\end{tabular}
\end{table}

Bảng~\ref{tab:det-results} cho thấy cả hai mô hình đạt mAP@0,5 trên 0,99 sau 50 epoch.
YOLOv8n có Precision cao hơn (0,988 so với 0,975), còn YOLO26n có mAP@0,5:0,95 cao hơn nhẹ
(0,881 so với 0,869); chênh lệch mAP@0,5 giữa hai mô hình dưới 0,002, nằm trong nhiễu thống
kê của ba seed. Đường cong huấn luyện theo epoch ở Hình~\ref{fig:det-curves}, hội tụ tốt từ
epoch 10 đến 15 với độ lệch chuẩn mAP@0,5 dưới 0,001 giữa các seed.

\begin{figure}[htbp]
\centering
\includegraphics[width=0.92\textwidth]{plate_det_curves.png}
\caption{Đường cong mAP theo epoch trên tập kiểm định, trung bình của ba seed (vùng bóng: độ lệch chuẩn). Trái: mAP@0,5; Phải: mAP@0,5:0,95.}

\label{fig:det-curves}
\end{figure}

Trên tập kiểm tra, YOLOv8n đạt mAP@0,5 = 0,989, mAP@0,5:0,95 = 0,864,
P = 0,990, R = 0,982; YOLO26n đạt mAP@0,5 = 0,980, mAP@0,5:0,95 = 0,859.

\subsection{Đánh Giá Sau Khử Trùng Lặp Gần}

Do khoảng 40\% ảnh kiểm tra của A1 có ảnh huấn luyện gần trùng (Mục~\ref{sec:dedup}),
YOLOv8n (seed 0) được đánh giá lại sau khi loại 230 ảnh này (Bảng~\ref{tab:det-dedup}).

\begin{table}[htbp]
\centering
\caption{Hiệu năng YOLOv8n (seed 0) trên tập kiểm tra trước và sau khi khử trùng lặp gần (Hamming $\leq 5$).}
\label{tab:det-dedup}
\begin{tabular}{lcccccc}
\toprule
\textbf{Tập kiểm tra} & \textbf{Số ảnh} & \textbf{P} & \textbf{R} & \textbf{mAP@0,5} & \textbf{mAP@0,5:0,95} \\
\midrule
Gốc (có near-dup) & 572 & 0,989 & 0,986 & 0,989 & 0,864 \\
Đã khử trùng lặp  & 342 & 0,983 & 0,980 & 0,983 & 0,834 \\
\midrule
\textit{Chênh lệch}  & \textit{$-$230} & \textit{$-$0,006} & \textit{$-$0,006} & \textit{$-$0,006} & \textit{$-$0,030} \\
\bottomrule
\end{tabular}
\end{table}

Sau khi loại ảnh gần trùng, mAP@0,5 chỉ giảm 0,006 (từ 0,989 xuống 0,983) nhưng
mAP@0,5:0,95 giảm 0,030 (từ 0,864 xuống 0,834): mô hình đã ghi nhớ vị trí biển trong các khung
gần trùng nên đạt IoU cao khi gặp lại. Vì vậy mAP@0,5 bằng 0,983 là ước lượng bảo thủ và đáng
tin hơn cho dữ liệu chưa thấy.

\section{Thảo Luận}
\label{sec:phathien-thaoluan}

YOLO26n nhỉnh hơn ở mAP@0,5:0,95, gợi ý việc bỏ DFL và dùng STAL cải thiện định vị; YOLOv8n
có Precision cao hơn, tức ít báo nhầm, có lợi ở bãi xe vì một phát hiện nhầm dẫn tới lỗi tra
biển. YOLOv8n chỉ khoảng 3,2 triệu tham số và 8,7 GFLOPs~\autocite{jocher2023yolov8}.

Đề tài chọn YOLOv8n (seed 0) cho pipeline vì hệ sinh thái xuất ONNX trên CPU ARM trưởng thành
hơn và Precision cao hơn, trong khi mAP@0,5 của hai mô hình tương đương. YOLO26n là ứng viên
nâng cấp khi tối ưu thiết bị biên (Mục~\ref{sec:huongphattrien}).
